% Preamble: \usepackage{graphicx}
% Copy the three PDFs into figures/ and use \graphicspath{{figures/}}.
% In a two-column manuscript, use figure* instead of figure to span both columns.

% Adjustable placement parameters:
\providecommand{\StageTwoPlotWidth}{\linewidth}
\providecommand{\StageTwoDistanceWidth}{0.44\linewidth}
\providecommand{\StageTwoPriorWidth}{0.47\linewidth}
\providecommand{\StageTwoLegendGap}{0.04\linewidth}
\providecommand{\StageTwoVerticalGap}{2mm}

\begin{figure}[tbp]
    \centering
    \begin{minipage}[t]{\StageTwoDistanceWidth}
        \centering
        \vspace{0pt}
        \includegraphics[width=\linewidth]{stage2_prior_ler_legend_distance.pdf}
    \end{minipage}%
    \hspace{\StageTwoLegendGap}%
    \begin{minipage}[t]{\StageTwoPriorWidth}
        \centering
        \vspace{0pt}
        \includegraphics[width=\linewidth]{stage2_prior_ler_legend_decoder_alias.pdf}
    \end{minipage}

    \vspace{\StageTwoVerticalGap}
    \includegraphics[width=\StageTwoPlotWidth]{stage2_prior_ler.pdf}
    \caption{Comparison of the stage-2 prior choices in prior-perturbation
    ensemble decoding for triangular color codes with $d=7,9,11$, one round,
    and independent bit-flip noise. Both variants use perturbed stage-1 priors;
    stage 2 uses either the original prior or the perturbed prior of the same
    ensemble member. The ensemble contains $M=16$ members, including the
    unperturbed baseline, with perturbation strength $\alpha=1$.
    Final candidates are selected using the original X/Z DEM log odds in
    both variants. Each point uses $10^6$ shots; shaded bands show 99\%
    Wilson confidence intervals. Colors identify the distance and markers
    identify the stage-2 prior choice.}
    \label{fig:stage2-prior-comparison}
\end{figure}
